\newpage
\phantomsection
\label{prf:computational-geometry-greedy-slope-upper-hull}
\LRAProofFor{thm:computational-geometry-greedy-slope-upper-hull}

\begin{remark*}[Return]
\hyperref[thm:computational-geometry-greedy-slope-upper-hull]{Return to Theorem}
\end{remark*}
\ProofVaultURL{proof-vault/pending/computational-geometry-greedy-slope-upper-hull}

\begin{theorem*}[Greedy slope construction of the upper hull]
Let \(P\subset\mathbb{R}^2\) be finite and in general position. Define
\[
  v_0=\argmin_{p\in P}x(p).
\]
Given \(v_k\) with
\[
  x(v_k)<\max_{p\in P}x(p),
\]
define
\[
  v_{k+1}
  =
  \argmax_{\substack{q\in P\\x(q)>x(v_k)}}
  \frac{y(q)-y(v_k)}{x(q)-x(v_k)}.
\]
Terminate once
\[
  v_k=\argmax_{p\in P}x(p).
\]
Then the points
\[
  v_0,v_1,\dots,v_m
\]
produced by this construction are exactly the vertices of the upper hull of
\(P\), and the slopes encountered along the way are strictly decreasing.
\end{theorem*}

\begin{proof}
\textbf{Professional Standard Proof.}~\\
TODO: Show that each greedy slope choice selects the next upper-hull vertex,
then use the monotone slope theorem to prove that no upper-hull vertex is
skipped.
\end{proof}

\begin{proof}
\textbf{Detailed Learning Proof.}~\\
TODO: Expand the induction on the produced vertices and identify where
general-position uniqueness is used.
\end{proof}

\begin{remark*}[Proof structure]
The proof is the algorithmic consequence of the monotone slope
characterization.
\end{remark*}

\begin{dependencies}
\begin{itemize}
  \item \hyperref[def:computational-geometry-convex-hull]{Convex Hull}
  \item \hyperref[thm:computational-geometry-monotone-slope-upper-hull]{Monotone Slope Characterization of the Upper Hull}
  \item TODO: finite set maximum
  \item TODO: slope of a line
\end{itemize}
\end{dependencies}

\clearpage
